\chapter{Chapter07 — CSP Examples}
%%% generated by the blueprint pipeline from FABL.Chapter07.CSPExamples
\section{Overview}

Declarations of the Lean module \texttt{FABL.Chapter07.CSPExamples}.

\section{Declarations}

\begin{definition}[Scope of a max cut constraint]
\lean{FABL.maxCutConstraintScope}
\statementsource{Analysis-of-Boolean-Functions-by-Ryan-ODonnell}{
  pdf-ffbe9c7b3306-p185-b004,
  pdf-ffbe9c7b3306-p185-b005
}
\label{FABL.maxCutConstraintScope}
\leanok
\uses{FABL.CSPConstraint, FABL.maxCutTemplate}
For a constraint in the max cut template on a variable set $V$, its scope is the
embedding of the two positions of the constraint into $V$; that is, it sends each
position to the variable assigned to that position by the constraint.
\end{definition}

\begin{theorem}[Distinct variables in a cut constraint]
\lean{FABL.maxCutConstraint_scope_ne}
\proofsource{Analysis-of-Boolean-Functions-by-Ryan-ODonnell}{
  pdf-ffbe9c7b3306-p185-b004,
  pdf-ffbe9c7b3306-p185-b005
}
\label{FABL.maxCutConstraint_scope_ne}
\leanok
\uses{FABL.CSPConstraint, FABL.maxCutConstraintScope, FABL.maxCutTemplate}
\difficulty{4}
For any set $V$ and any maximum-cut constraint on $V$, the two variables in its scope
are distinct.
\end{theorem}

\begin{theorem}[Max cut instance exists incident constraint]
\lean{FABL.maxCutInstance_exists_incident_constraint}
\statementsource{Analysis-of-Boolean-Functions-by-Ryan-ODonnell}{pdf-ffbe9c7b3306-p185-b004}
\label{FABL.maxCutInstance_exists_incident_constraint}
\leanok
\uses{FABL.CSPInstance, FABL.maxCutConstraintScope, FABL.maxCutTemplate}
\difficulty{4}
% dataset-pipeline: stub -- prose pending build_informal_statements.py
Statement pending regeneration.
\end{theorem}

\begin{definition}[Two vertex max cut constraint]
\lean{FABL.twoVertexMaxCutConstraint}
\statementsource{Analysis-of-Boolean-Functions-by-Ryan-ODonnell}{pdf-ffbe9c7b3306-p184-b003}
\label{FABL.twoVertexMaxCutConstraint}
\leanok
\uses{FABL.CSPConstraint, FABL.maxCutTemplate}
On the vertex set $\{0,1\}$, the two vertex max cut constraint uses the unique max cut
predicate and has scope equal to the identity map. It is satisfied exactly when the
labels at the two vertices differ.
\end{definition}

\begin{definition}[Parallel edge max cut instance]
\lean{FABL.parallelEdgeMaxCutInstance}
\statementsource{Analysis-of-Boolean-Functions-by-Ryan-ODonnell}{
  pdf-ffbe9c7b3306-p185-b004,
  pdf-ffbe9c7b3306-p185-b005
}
\label{FABL.parallelEdgeMaxCutInstance}
\leanok
\uses{FABL.CSPInstance, FABL.maxCutTemplate, FABL.twoVertexMaxCutConstraint}
On the vertex set $\{0,1\}$, take the max cut instance consisting of two copies of the
constraint joining the two vertices. The constraint multiset is nonempty, and each
vertex occurs in the scope of a constraint.
\end{definition}

\begin{theorem}[Constraint count of the parallel-edge max-cut instance]
\lean{FABL.parallelEdgeMaxCutInstance_card}
\label{FABL.parallelEdgeMaxCutInstance_card}
\leanok
\uses{FABL.parallelEdgeMaxCutInstance}
\difficulty{1}
The parallel-edge max-cut instance has exactly two constraints.
\end{theorem}

\begin{definition}[Constant assignment over the two element field]
\lean{FABL.e3LinConstantAssignment}
\statementsource{Analysis-of-Boolean-Functions-by-Ryan-ODonnell}{pdf-ffbe9c7b3306-p188-b005}
\label{FABL.e3LinConstantAssignment}
\leanok
\uses{FABL.𝔽₂}
For a variable set $V$ and a label $b$ in the field with two elements, the constant
assignment with label $b$ is the function on $V$ that assigns $b$ to every variable.
\end{definition}

\begin{theorem}[Constant Assignment Satisfaction Indicator]
\lean{FABL.satisfiedIndicator_e3LinConstantAssignment}
\proofsource{Analysis-of-Boolean-Functions-by-Ryan-ODonnell}{
  pdf-ffbe9c7b3306-p188-b005,
  pdf-ffbe9c7b3306-p185-b002
}
\label{FABL.satisfiedIndicator_e3LinConstantAssignment}
\leanok
\uses{FABL.CSPConstraint, FABL.CSPConstraint.satisfiedIndicator, FABL.e3LinConstantAssignment, FABL.maxE3LinTemplate, FABL.𝔽₂}
\difficulty{5}
Let $c$ be a constraint over $\mathbb{F}_2$ that requires the sum of the values on its
three pairwise-distinct variables to equal its right-hand side $b$. Under the assignment
giving every variable the same value $b$, the satisfaction indicator of $c$ is $1$ if
the constraint's right-hand side is $b$, and $0$ otherwise.
\end{theorem}

\begin{theorem}[Complementary Satisfaction Indicators]
\lean{FABL.satisfiedIndicator_zero_add_one_e3Lin}
\proofsource{Analysis-of-Boolean-Functions-by-Ryan-ODonnell}{pdf-ffbe9c7b3306-p185-b002}
\label{FABL.satisfiedIndicator_zero_add_one_e3Lin}
\leanok
\uses{FABL.CSPConstraint, FABL.CSPConstraint.satisfiedIndicator, FABL.boolF₂Equiv, FABL.e3LinConstantAssignment, FABL.maxE3LinTemplate, FABL.satisfiedIndicator_e3LinConstantAssignment, FABL.𝔽₂}
\difficulty{4}
Let $c$ be a constraint over the field with two elements, whose predicate requires the
sum of its three scoped values to equal a specified right-hand side. The satisfaction
indicators of $c$ under the two constant assignments taking the values $0$ and $1$ sum
to $1$.
\end{theorem}

\begin{theorem}[Complementary satisfaction counts for constant assignments]
\lean{FABL.sum_satisfiedIndicators_zero_add_one_e3Lin}
\label{FABL.sum_satisfiedIndicators_zero_add_one_e3Lin}
\leanok
\uses{FABL.CSPConstraint, FABL.CSPConstraint.satisfiedIndicator, FABL.e3LinConstantAssignment, FABL.maxE3LinTemplate, FABL.satisfiedIndicator_zero_add_one_e3Lin}
\difficulty{4}
Let $V$ be any set, and let $\mathcal C$ be a finite multiset of constraints, each
requiring that the sum in $\mathbb F_2$ of the values on three distinct variables equal
a specified right-hand side. The number of constraints satisfied by the assignment that
gives every variable value $0$, plus the number satisfied by the assignment that gives
every variable value $1$, equals the total number of constraints in $\mathcal C$.
\end{theorem}

\begin{theorem}[Complementary Constant Assignments]
\lean{FABL.satisfiedCount_zero_add_one_e3Lin}
\label{FABL.satisfiedCount_zero_add_one_e3Lin}
\leanok
\uses{FABL.CSPInstance, FABL.CSPInstance.satisfiedCount, FABL.e3LinConstantAssignment, FABL.maxE3LinTemplate, FABL.sum_satisfiedIndicators_zero_add_one_e3Lin}
\difficulty{1}
Let $P$ be a finite constraint instance over $\mathbb{F}_2$ whose constraints are
equations on three pairwise-distinct variables, with at least one constraint and every
variable occurring in some constraint. For each constraint, satisfaction means that the
sum of its three assigned values equals its right-hand side. The total numbers of
satisfied constraints under the assignments that give every variable $0$ and every
variable $1$ sum to the total number of constraints, counted with multiplicity.
\end{theorem}

\begin{theorem}[Constant assignment for three-variable linear constraints]
\lean{FABL.exists_constantAssignment_value_ge_half_e3Lin}
\proofsource{Analysis-of-Boolean-Functions-by-Ryan-ODonnell}{pdf-ffbe9c7b3306-p188-b005}
\label{FABL.exists_constantAssignment_value_ge_half_e3Lin}
\leanok
\uses{FABL.CSPInstance, FABL.CSPInstance.satisfiedCount, FABL.CSPInstance.value, FABL.e3LinConstantAssignment, FABL.maxE3LinTemplate, FABL.satisfiedCount_zero_add_one_e3Lin, FABL.𝔽₂}
\difficulty{4}
Let $P$ be a finite instance of the constraint problem over $\mathbb F_2$ whose
predicates require the sum of three distinct variables to equal a specified right-hand
side. Then one of the two constant assignments, assigning the same bit to every
variable, satisfies at least half of the constraints, counted with multiplicity.
\end{theorem}

\begin{definition}[Better constant label for three variable equations]
\lean{FABL.betterE3LinConstantLabel}
\statementsource{Analysis-of-Boolean-Functions-by-Ryan-ODonnell}{pdf-ffbe9c7b3306-p188-b005}
\label{FABL.betterE3LinConstantLabel}
\leanok
\uses{FABL.CSPInstance, FABL.CSPInstance.value, FABL.e3LinConstantAssignment, FABL.maxE3LinTemplate, FABL.𝔽₂}
Let $P$ be an instance of equations over the field with two elements, each involving
three variables. Compare its values on the two constant assignments: the selected label
is $1$ if the value at the all-zero assignment is at most the value at the all-one
assignment, and is $0$ otherwise.
\end{definition}

\begin{definition}[Better constant assignment for three variable equations]
\lean{FABL.betterE3LinConstantAssignment}
\statementsource{Analysis-of-Boolean-Functions-by-Ryan-ODonnell}{pdf-ffbe9c7b3306-p188-b005}
\label{FABL.betterE3LinConstantAssignment}
\leanok
\uses{FABL.CSPInstance, FABL.betterE3LinConstantLabel, FABL.e3LinConstantAssignment, FABL.maxE3LinTemplate, FABL.𝔽₂}
Let $V$ be a type and let $P$ be a nonempty finite multiset of constraints from the
template of equations $x_1+x_2+x_3=b$ over $\bbr/	ext{?}$. Define the better constant
assignment of $P$ by assigning each variable the value selected by the construction
`betterE3LinConstantLabel` followed by `e3LinConstantAssignment`.
\end{definition}

\begin{theorem}[Value of a constant assignment for three-variable equations]
\lean{FABL.betterE3LinConstantAssignment_value_ge_half}
\proofsource{Analysis-of-Boolean-Functions-by-Ryan-ODonnell}{pdf-ffbe9c7b3306-p188-b005}
\label{FABL.betterE3LinConstantAssignment_value_ge_half}
\leanok
\uses{FABL.CSPInstance, FABL.CSPInstance.value, FABL.betterE3LinConstantAssignment, FABL.boolF₂Equiv, FABL.e3LinConstantAssignment, FABL.exists_constantAssignment_value_ge_half_e3Lin, FABL.maxE3LinTemplate}
\difficulty{5}
Let $V$ be a type of variables and consider a finite multiset of constraints, with at
least one constraint, in which every variable occurs in some constraint. Each constraint
selects an ordered triple of distinct variables and requires their sum in the field with
two elements to equal a specified element of that field. Choose the constant assignment,
giving every variable the same field element, whose instance value is larger (choosing
the assignment with value $1$ in case of a tie). Its fraction of satisfied constraints
is at least $1/2$.
\end{theorem}

\begin{theorem}[Lower bound on the linear constraint optimum]
\lean{FABL.half_le_optimum_maxE3Lin}
\proofsource{Analysis-of-Boolean-Functions-by-Ryan-ODonnell}{pdf-ffbe9c7b3306-p188-b005}
\label{FABL.half_le_optimum_maxE3Lin}
\leanok
\uses{FABL.CSPInstance, FABL.CSPInstance.optimum, FABL.CSPInstance.value_le_optimum, FABL.betterE3LinConstantAssignment, FABL.betterE3LinConstantAssignment_value_ge_half, FABL.maxE3LinTemplate}
\difficulty{1}
Let $V$ be a finite set and let $P$ be an instance of the three-variable linear
constraint problem over the field with two elements, with a nonempty multiset of
constraints and every variable in $V$ occurring in some constraint scope. Then the
maximum fraction of constraints satisfied by an assignment of field elements to $V$ is
at least $1/2$.
\end{theorem}
